\chapter{Introduction}
\label{ch:introduction}

% ~6-10 pages

\section{Clinical motivation}
% Outline:
% - DR prevalence (mention WHO numbers, global blindness statistics)
% - 5-stage severity scale and its treatment implications:
%   * No DR / Mild — observation
%   * Moderate — closer follow-up
%   * Severe / PDR — urgent referral, possible laser/anti-VEGF
% - Cost and shortage of trained ophthalmologists, especially in
%   primary care; need for automated triage.

\section{Why uncertainty matters in medical AI}
% Outline:
% - Modern CNNs are over-confident: a wrong 99\%-confident prediction is
%   indistinguishable from a right one.
% - Risk of harm: an overconfident "No DR" can delay vision-saving care.
% - Recent regulatory push (FDA AI-as-a-medical-device, EU AI Act)
%   explicitly requires transparency around model uncertainty.
% - Clinically useful systems should produce a "refer to specialist"
%   recommendation rather than a forced binary verdict on hard cases.

\section{Bachelor recap and research gap}
% Outline:
% - Briefly describe the bachelor thesis (binary classification, 9
%   algorithms, single 80/20 split, accuracy as the only metric).
% - Identify the gaps the master must close:
%   * Multi-class severity grading, not just binary
%   * Statistical significance and confidence intervals
%   * Calibration, not just accuracy
%   * Distribution-free uncertainty guarantees
%   * Out-of-distribution detection
%   * Reproducibility under cross-validation

\section{Contributions}
% Outline:
% 1. A reproducible, statistically rigorous evaluation pipeline for the
%    APTOS 2019 dataset (stratified splits, K-fold CV, McNemar testing).
% 2. Calibration analysis (ECE, MCE, temperature scaling) for six DL
%    architectures, exposing previously unreported over-confidence.
% 3. Application of split conformal prediction (LAC + APS) to DR
%    classification with formal $1-\alpha$ coverage guarantees.
% 4. Comparison of Monte Carlo Dropout, deep-ensemble, and feature-space
%    OOD detection methods on the same test set.
% 5. A 5-stage grading model achieving $\text{QWK} = 0.847$, competitive
%    with the APTOS 2019 Kaggle leaderboard.
% 6. An open-source Streamlit application that surfaces uncertainty
%    visually for end users.

\section{Thesis structure}
% \cref{ch:background} surveys clinical and technical background;
% \cref{ch:methodology} describes the data, splits and shared experimental
% protocol; \cref{ch:phase1} reports binary baselines and calibration;
% \cref{ch:phase2} introduces conformal prediction, OOD detection,
% MC Dropout and K-fold CV; \cref{ch:phase3} extends the system to
% multi-class severity grading; \cref{ch:discussion} discusses clinical
% implications, limitations and future work.
